\input{../tex-inputs/latex-leseansicht-vorspann}

\section[Arthur Schnitzler an Gustav Schwarzkopf, 6. 7. 1914]{L04491 Arthur Schnitzler an Gustav Schwarzkopf, 6. 7. 1914}
\nopagebreak\mylabel{L04491v}
\rehead{ }\normalsize\beginnumbering\briefempfaengerindex{Schwarzkopf, Gustav@\textsc{Schwarzkopf, Gustav}!zzzSchnitzler, Arthur@\emph{von Arthur Schnitzler}!1914-07-061@{6. 7. 1914}|(be}
\toendnotes[C]{\smallbreak\pagebreak[2]}
\correspDesc{Versand  durch Arthur Schnitzler am 6. 7. 1914 in Wien
\newline{}Übermittlung  am 6. 7. 1914 in Wien
\newline{}Erhalt  durch Gustav Schwarzkopf im Zeitraum [7. 7. 1914
                  – 11. 7. 1914?] in Weißenbach am Attersee}\toendnotes[C]{\smallbreak}
\Standort{DLA, A:Schnitzler, HS.1985.1.1897.}
\physDesc{Bildpostkarte, 270 Zeichen
\newline{}Handschrift: schwarze Tinte, deutsche Kurrent
\newline{}Versand: Stempel: »\nobreak{}\oindex{Wien@\textbf{Wien}|pwk}1/1 Wien 8, 6. VII. 14, 7\nobreak{}«.  }\toendnotes[C]{\smallbreak}\pstart{}Herrn Guſtav Schwarzkopf\pend{}\pstart{}Weißenbach\oindex{Weißenbach am Attersee@\textbf{Weißenbach am Attersee}|pw}\pend{}\pstart{}am Atterſee\oindex{Schloss Leopoldskron@\textbf{Schloss Leopoldskron}|pw}\pend{}\pstart{}(pr. Adr. Dr. \textsc{Beer Hofman\pwindex{Beer-Hofmann, Richard 11.\,7.\,1866 Wien – 26.\,9.\,1945 New York City@\textsc{Beer-Hofmann, Richard} (11.\,7.\,1866 Wien – 26.\,9.\,1945 New York City)|pw}})\pend{}{\bigskip}
\pstart
           \noindent{}{\pb}\textcolor{gray}{\textbf{Wien, XVIII, Sternwartestr. 71\oindex{Wien@\textbf{Wien}!XVIII., Währing@\textbf{XVIII., Währing}!Sternwartestraße 71@\textbf{Sternwartestraße 71}|pw}.}}\pend
           \vspace{1em}
\pstart
           \noindent{}Herzlichen Gruſs u Dank für die \label{K_L04491-1v}\edtext{Karte}{\lemma{\textnormal{\emph{Karte}}}\Cendnote{\textnormal{XXXX Auszeichnungsfehler: Dokument L04249 nicht gefunden.}}}\label{K_L04491-1}! Wir wollten
                  Samſtag für ein paar Tage \label{K_L04491-2v}\edtext{nach \uline{Lunz}\oindex{Lunz am See@\textbf{Lunz am See}|pw}}{\lemma{\textnormal{\emph{nach Lunz}}}\Cendnote{\textnormal{Erst am 14. 7. 1914 brachen Arthur, Olga\pwindex{Schnitzler, Olga 17.\,1.\,1882 Wien – 13.\,1.\,1970 Lugano@\textsc{Schnitzler, Olga} (17.\,1.\,1882 Wien – 13.\,1.\,1970 Lugano)|pwk} und Heinrich Schnitzler\pwindex{Schnitzler, Heinrich 9.\,8.\,1902 Hinterbrühl – 12.\,7.\,1982 Wien@\textsc{Schnitzler, Heinrich} (9.\,8.\,1902 Hinterbrühl – 12.\,7.\,1982 Wien)|pwk} nach Lunz am See\oindex{Lunz am See@\textbf{Lunz am See}|pwk} auf.}}}\label{K_L04491-2}, (und 8 Tage){ }ſpäter nach \label{K_L04491-3v}\edtext{Pontreſina\oindex{Pontresina@\textbf{Pontresina}|pw}}{\lemma{\textnormal{\emph{Pontresina}}}\Cendnote{\textnormal{Nach einem Aufenthalt in Thusis\oindex{Thusis@\textbf{Thusis}|pwk} ab dem 18. 7. 1914 gelangte die Familie Schnitzler\pwindex{Schnitzler, Olga 17.\,1.\,1882 Wien – 13.\,1.\,1970 Lugano@\textsc{Schnitzler, Olga} (17.\,1.\,1882 Wien – 13.\,1.\,1970 Lugano)|pwk}\pwindex{Schnitzler, Heinrich 9.\,8.\,1902 Hinterbrühl – 12.\,7.\,1982 Wien@\textsc{Schnitzler, Heinrich} (9.\,8.\,1902 Hinterbrühl – 12.\,7.\,1982 Wien)|pwk}\pwindex{Cappellini, Lili 13.\,9.\,1909 Wien – 26.\,7.\,1928 Venedig@\textsc{Cappellini, Lili} (13.\,9.\,1909 Wien – 26.\,7.\,1928 Venedig)|pwk} am
                     21. 7. 1914 nach Pontresina\oindex{Pontresina@\textbf{Pontresina}|pwk} und zwei Tag darauf nach Celerina\oindex{Celerina@\textbf{Celerina}|pwk}, wo sie gut zwei Wochen
                  verbrachte.}}}\label{K_L04491-3}.\pend
           
\pstart
           Wir hoffen Sie und \label{K_L04491-4v}\edtext{Richard’s\pwindex{Beer-Hofmann, Richard 11.\,7.\,1866 Wien – 26.\,9.\,1945 New York City@\textsc{Beer-Hofmann, Richard} (11.\,7.\,1866 Wien – 26.\,9.\,1945 New York City)|pw}\pwindex{Beer-Hofmann, Paula 25.\,2.\,1879 Wien – 30.\,10.\,1939 Zürich@\textsc{Beer-Hofmann, Paula} (25.\,2.\,1879 Wien – 30.\,10.\,1939 Zürich)|pw}\pwindex{Beer-Hofmann, Gabriel 9.\,1.\,1901 Wien – 24.\,3.\,1971 St Albans@\textsc{Beer-Hofmann, Gabriel} (9.\,1.\,1901 Wien – 24.\,3.\,1971 St Albans)|pw}\pwindex{Beer-Hofmann, Naëmah 20.\,12.\,1898 Wien – 10.\,11.\,1971 New York City@\textsc{Beer-Hofmann, Naëmah} (20.\,12.\,1898 Wien – 10.\,11.\,1971 New York City)|pw}\pwindex{Beer-Hofmann, Mirjam 4.\,9.\,1897 Wien – 24.\,12.\,1984 New York City@\textsc{Beer-Hofmann, Mirjam} (4.\,9.\,1897 Wien – 24.\,12.\,1984 New York City)|pw}}{\lemma{\textnormal{\emph{Richard’s}}}\Cendnote{\textnormal{Die Familie Beer-Hofmann\pwindex{Beer-Hofmann, Richard 11.\,7.\,1866 Wien – 26.\,9.\,1945 New York City@\textsc{Beer-Hofmann, Richard} (11.\,7.\,1866 Wien – 26.\,9.\,1945 New York City)|pwk}\pwindex{Beer-Hofmann, Paula 25.\,2.\,1879 Wien – 30.\,10.\,1939 Zürich@\textsc{Beer-Hofmann, Paula} (25.\,2.\,1879 Wien – 30.\,10.\,1939 Zürich)|pwk}\pwindex{Beer-Hofmann, Gabriel 9.\,1.\,1901 Wien – 24.\,3.\,1971 St Albans@\textsc{Beer-Hofmann, Gabriel} (9.\,1.\,1901 Wien – 24.\,3.\,1971 St Albans)|pwk}\pwindex{Beer-Hofmann, Naëmah 20.\,12.\,1898 Wien – 10.\,11.\,1971 New York City@\textsc{Beer-Hofmann, Naëmah} (20.\,12.\,1898 Wien – 10.\,11.\,1971 New York City)|pwk}\pwindex{Beer-Hofmann, Mirjam 4.\,9.\,1897 Wien – 24.\,12.\,1984 New York City@\textsc{Beer-Hofmann, Mirjam} (4.\,9.\,1897 Wien – 24.\,12.\,1984 New York City)|pwk} verbrachte den
                  Sommer 1914 in Weißenbach am
                  Attersee\oindex{Weißenbach am Attersee@\textbf{Weißenbach am Attersee}|pwk}.}}}\label{K_L04491-4} wohl u vergnügt. Laſſen Sie öfters von ſich hören.\pend
           \pstart Ihr \spacefill\mbox{A}.\pend{}\selectlanguage{ngerman}\endnumbering\briefempfaengerindex{Schwarzkopf, Gustav@\textsc{Schwarzkopf, Gustav}!zzzSchnitzler, Arthur@\emph{von Arthur Schnitzler}!1914-07-061@{6. 7. 1914}|)be}\mylabel{L04491h}
\begin{anhang}
\end{anhang}\newcommand{\dateiname}{L04491}\newcommand{\titel}{Arthur Schnitzler an Gustav Schwarzkopf, 6. 7. 1914}\newcommand{\editorInnen}{Herausgegeben von Jahnke, SelmaMüller, Martin Anton}\input{../tex-inputs/latex-leseansicht-abspann}
